\ncertchapter{12}{Collisions}\chlabel{ch:collisions}

A collision is a brief and violent interaction: two bodies meet, exert enormous forces on each other for a
very short time, and separate, or stick together, with altered velocities. The forces during the impact are
almost never known in detail and are not needed. What is needed is a way to connect the velocities before
the collision to the velocities after it without knowing what happened in between, and the conservation laws
provide exactly that.

The entire subtlety of the chapter lies in knowing which conservation law always holds and which does
not. Momentum is conserved in every collision between free bodies; kinetic energy only in the special
case of an elastic collision, and assuming otherwise is the most common error in this part of the
syllabus. The chapter first proves the momentum claim, then works through the possible fates of the
kinetic energy: fully kept (elastic), most lost (the bodies stick), and everything in between, measured
by a coefficient of restitution. It closes with glancing collisions in two dimensions.

\section{Why Momentum Is Always Conserved}

The claim that momentum is conserved in every collision between free bodies, whatever the materials and
however violent the impact, needs only two ingredients: the third law, and the brevity of the collision.
(\chref{ch:newton} met the same argument; here we see exactly where each ingredient enters.)

\begin{derivbox}{Momentum conservation in a collision}
Two bodies collide. Body~1 has momentum $\vp_1=m_1\vv_1$ and body~2 has $\vp_2=m_2\vv_2$ (momentum, a
vector in \si{kg.m/s}, from \chref{ch:newton}). During contact, body~1 pushes on body~2 with a force
$\vF_{12}$, and body~2 pushes back on body~1 with $\vF_{21}$.

\stepn{1} \textbf{Second law for each body, third law between them.} Force is the rate of change of
momentum, so $d\vp_1/dt=\vF_{21}$ and $d\vp_2/dt=\vF_{12}$, while the third law makes
$\vF_{21}=-\vF_{12}$ at every instant.

\stepn{2} \textbf{Add.} The total momentum $\vc{P}=\vp_1+\vp_2$ changes at the rate
\[
\frac{d\vc{P}}{dt}=\vF_{21}+\vF_{12}=\vc0,
\]
however large, brief, or complicated the contact forces are. With subscripts $i$ and $f$ meaning just
before and just after the collision,
\[
\boxed{\;\vp_{1i}+\vp_{2i}=\vp_{1f}+\vp_{2f}\;}
\]
You never need to know anything about the messy forces during contact.
\end{derivbox}

The derivation assumed that no external force acts, and the brevity of the collision is what makes that
nearly true. Recall from \chref{ch:newton} that a force acting for a time $\Delta t$ changes momentum by its
impulse, $\vc{J}=\int\vF\,dt$ (a vector, in \si{N.s}), about $\vF_{\text{avg}}\,\Delta t$ where
$\vF_{\text{avg}}$ is the average force. A \SI{0.16}{kg} cricket ball meeting a bat at
\SI{30}{m/s} and leaving at \SI{30}{m/s} the other way changes its momentum by about \SI{9.6}{N.s} in a
contact of about \SI{1}{ms}; in that millisecond gravity supplies an impulse of only
$mg\,\Delta t\approx\SI{0.0016}{N.s}$, less than a five-thousandth as much. Any external force of
ordinary size (gravity, friction) is negligible during a collision in this way. Forces that are
themselves \textbf{impulsive}, huge and brief like the contact force, are not: the normal force of a
floor, or the force at a hinge, can deliver an impulse as large as the collision's own. So the
momentum of the colliding bodies alone is conserved only when neither is held by such a support.

\section{Classifying Collisions}

Momentum conservation is one equation, but a head-on collision with known masses and initial velocities
has two unknowns, the two final velocities. A second piece of information is needed, and what it can be
depends on what happens to the kinetic energy. That is why collisions are classified by their kinetic
energy, and not by their momentum, which is conserved in all of them (Figure~\ref{fig:collasym}).

\begin{keybox}{What is conserved, and when}
\textbf{The total momentum of the colliding bodies is conserved in every collision in which external
impulses are negligible}; this fails when one body is fixed or pivoted (a wall, a floor, a hinged door).
\textbf{Kinetic energy is conserved
only in an \emph{elastic} collision.} In an \emph{inelastic} one, some kinetic energy becomes heat,
sound, and permanent deformation; in a \emph{perfectly inelastic} collision the bodies stick together
and the kinetic-energy loss is the greatest that momentum conservation allows. The ``lost'' kinetic
energy is not destroyed; it becomes internal energy, so \emph{total} energy is conserved in every case.
(If stored energy is released at contact, as by a compressed spring or an explosive cap, kinetic energy
can even \emph{increase}: a superelastic collision.)
\end{keybox}

\begin{figure}[htbp]
\centering
\tikzfig{collision_asymmetry}
\figcaption{fig:collasym}{The asymmetry at the heart of the chapter. Momentum is protected because the internal forces act in equal and opposite pairs (left). Kinetic energy has no such protection (right): some of it may go into deformation, heat and sound, and the sound of the impact is a sign of that loss.}
\end{figure}

Between the two extremes the second equation is supplied by one measured number, which also labels the
kinds of collision.

\begin{defbox}{Coefficient of restitution}
The \textbf{line of impact} is the line along which the contact force acts: the common normal to the two
surfaces at the point of contact, which for two spheres is the line joining their centres. Let $u_1,u_2$
be the velocity components of the two bodies along this line just before impact and $v_1,v_2$ just after,
all \emph{signed}, with the positive direction pointing from body~1 towards body~2 (so before impact
$u_1>u_2$: body~1 is catching up). The
coefficient of restitution $e$ (a pure number) measures the ``bounciness'':
\[
e=\frac{\text{speed of separation}}{\text{speed of approach}}=\frac{v_2-v_1}{u_1-u_2}.
\]
$e=1$ is
perfectly elastic; $e=0$ is perfectly inelastic (they move together); $0<e<1$ is the real, partially
elastic case. The value of $e$ is an empirical property of the pair of surfaces (Newton's experimental
law): roughly constant, though in practice it depends somewhat on the impact speed.
\end{defbox}

\begin{trapbox}{put signs into $e$}
Use signed velocities, not speeds. A ball at \SI{4}{m/s} hitting a block at rest and bouncing back at
\SI{1}{m/s} while the block moves off at \SI{2}{m/s} has $v_1=-1$, so $e=[2-(-1)]/(4-0)=0.75$, not
$(2-1)/4=0.25$.
\end{trapbox}

\begin{remarkbox}{What happens during the impact}
Why does $e$ lie between 0 and 1? A collision has two phases.
 In the \emph{compression} phase the
bodies squash until, for an instant, they move with a common velocity $c$; in the \emph{restitution}
phase the deformation partly springs back and pushes them apart. Let the contact force deliver an
impulse $J_C$ (a signed component along the line of impact, in \si{N.s}) to body~2 during compression and
$J_R$ during restitution; body~1 receives the opposite impulses. Impulse equals change of momentum, so for
compression $m_2(c-u_2)=J_C$ and $m_1(c-u_1)=-J_C$. Dividing each by its mass and subtracting eliminates
$c$: $u_1-u_2=J_C\big(\tfrac1{m_1}+\tfrac1{m_2}\big)$. The same steps for restitution give
$v_2-v_1=J_R\big(\tfrac1{m_1}+\tfrac1{m_2}\big)$, so
$e=J_R/J_C$. A perfectly elastic pair springs back fully ($J_R=J_C$); putty does not spring back at all
($J_R=0$). This is why $e$ belongs to the pair of materials and lies between 0 and 1.
\end{remarkbox}

\section{Elastic Collisions in One Dimension}

Begin with the case in which kinetic energy itself supplies the second equation. In a \textbf{head-on}
(one-dimensional) collision all velocities lie along the line of impact, so we can drop the vector arrows
and work with the signed scalars $u_1,u_2,v_1,v_2$ defined above. (Momentum is still a vector law; in 1-D
it just becomes a statement about these signed numbers.)

\begin{derivbox}{One-dimensional elastic collision}
Masses $m_1,m_2$ with initial velocities $u_1,u_2$ collide head-on and elastically, leaving with
velocities $v_1,v_2$. ``Elastic'' gives us two conservation laws to work with:
\[
\underbrace{m_1u_1+m_2u_2=m_1v_1+m_2v_2}_{\text{momentum}},\qquad
\underbrace{\tfrac12m_1u_1^2+\tfrac12m_2u_2^2=\tfrac12m_1v_1^2+\tfrac12m_2v_2^2}_{\text{kinetic energy}}.
\]
\stepn{1} \textbf{Collect each mass on its own side.} Move the $m_1$ terms left and the $m_2$ terms right.
Momentum becomes $m_1(u_1-v_1)=m_2(v_2-u_2)$; energy (after cancelling the $\tfrac12$) becomes
$m_1(u_1^2-v_1^2)=m_2(v_2^2-u_2^2)$.
\stepn{2} \textbf{Factor the energy equation.} Each squared bracket is a difference of two squares:
$u_1^2-v_1^2=(u_1-v_1)(u_1+v_1)$, and likewise on the right. So the energy equation reads
\[
m_1(u_1-v_1)(u_1+v_1)=m_2(v_2-u_2)(v_2+u_2).
\]
\stepn{3} \textbf{Divide the energy equation by the momentum equation} (allowed because $u_1\ne v_1$: if
they were equal, no collision would have occurred). The factors
$m_1(u_1-v_1)$ and $m_2(v_2-u_2)$ are \emph{equal} (that is the momentum equation), so they cancel,
leaving just the ``sum'' factors:
\[
u_1+v_1=u_2+v_2\quad\Longrightarrow\quad \boxed{u_1-u_2=-(v_1-v_2)}.
\]
In words: the \emph{relative} velocity is unchanged in size but reversed in direction: the bodies
separate exactly as fast as they approached. (This is the meaning of $e=1$.)
\stepn{4} \textbf{Solve the two simple equations.} Step 3 gives $v_2=u_1+v_1-u_2$, which is linear.
Substituting it into momentum, $m_1u_1+m_2u_2=m_1v_1+m_2(u_1+v_1-u_2)$, so
$(m_1+m_2)v_1=(m_1-m_2)u_1+2m_2u_2$; then $v_2$ follows from Step 3:
\[
v_1=\frac{m_1-m_2}{m_1+m_2}\,u_1+\frac{2m_2}{m_1+m_2}\,u_2,\qquad
v_2=\frac{m_2-m_1}{m_1+m_2}\,u_2+\frac{2m_1}{m_1+m_2}\,u_1.
\]
Notice the symmetry: swapping the labels $1\leftrightarrow2$ turns one formula into the other, as it must.
\end{derivbox}

The formulas are easiest to remember, and to check, in their limiting cases (Figure~\ref{fig:collcases}).

\begin{figure}[htbp]
\centering
\tikzfig{collision_cases}
\figcaption{fig:collcases}{Three elastic head-on collisions with the target initially at rest, velocity arrows to scale. \emph{Equal masses:} the velocities swap (Newton's cradle, the billiard stop-shot). \emph{Heavy hits light} (here $m_1=4m_2$): the heavy body continues, only partly slowed, and knocks the light one ahead faster than it came. \emph{Light hits heavy} (here $m_2=4m_1$): the light body bounces back and the heavy one barely moves.}
\end{figure}

\begin{keybox}{The three limiting cases (target at rest, $u_2=0$)}
Put $u_2=0$ in the formulas above.
\textbf{Equal masses} ($m_1=m_2$): $v_1=0,\ v_2=u_1$, a complete transfer. \textbf{Heavy hits light}
($m_1\gg m_2$): $v_1\approx u_1,\ v_2\approx 2u_1$. \textbf{Light hits heavy} ($m_1\ll m_2$):
$v_1\approx-u_1,\ v_2\approx0$, the light body rebounds. The equal-mass transfer is why water, rich in
hydrogen, slows the fast neutrons in a nuclear reactor: a neutron has almost a proton's mass, so a
head-on hit on a hydrogen nucleus can stop it dead.
\end{keybox}

The general formulas still hide a simpler pattern, which makes them quicker to use. Let
$V=\dfrac{m_1u_1+m_2u_2}{m_1+m_2}$, the total momentum divided by the total mass; it is the velocity of
the pair's centre of mass (\chref{ch:com}), and since momentum is conserved it is the same after the
collision as before. The elastic formulas can be rewritten as
\[
v_1=2V-u_1,\qquad v_2=2V-u_2
\]
(check: $2V-u_1$ expands to the expression for $v_1$ above). Rearranged, $v_1-V=-(u_1-V)$: seen from a
frame moving with $V$, each body simply reverses its velocity.

\begin{exbox}{An elastic collision by the shortcut}
A \SI{2}{kg} ball at \SI{6}{m/s} strikes a \SI{1}{kg} ball at rest, head-on and elastically. Find both
final velocities.
\stepn{1} \textbf{Centre-of-mass velocity.} $V=\dfrac{2\times6+1\times0}{3}=\SI{4}{m/s}$.
\stepn{2} \textbf{Reverse in that frame.} $v_1=2(4)-6=\SI{2}{m/s}$ and $v_2=2(4)-0=\SI{8}{m/s}$.
\stepn{3} \textbf{Check.} Momentum: $2(2)+1(8)=12=2(6)$. Kinetic energy: $\tfrac12(2)(2^2)+\tfrac12(1)(8^2)
=4+32=\SI{36}{J}=\tfrac12(2)(6^2)$.\ \checkmark
\ans{$v_1=\SI{2}{m/s}$, $v_2=\SI{8}{m/s}$}
\end{exbox}

\begin{trapbox}{never assume kinetic energy is conserved}
Use kinetic-energy conservation only when the problem says elastic. Otherwise use momentum, and then
\emph{compute} the energy lost; do not set it to zero.
\end{trapbox}

\section{Perfectly Inelastic Collisions}\label{sec:inelastic12}

At the opposite extreme from the elastic case, the colliding bodies stick together ($e=0$). Now the
second equation is simply $v_1=v_2$, so the outcome is fixed by momentum alone. The questions are how
much kinetic energy is lost, and why this case loses the most.

\begin{derivbox}{Common velocity and energy lost}
Two bodies stick and move off together at a common velocity $v$.
\stepn{1} \textbf{Momentum (the only conservation law available).}
\[
m_1u_1+m_2u_2=(m_1+m_2)v\quad\Longrightarrow\quad
\boxed{\,v=\dfrac{m_1u_1+m_2u_2}{m_1+m_2}\,}.
\]
\stepn{2} \textbf{Kinetic energy lost.} Let $K_{\text{lost}}=K_i-K_f$, the initial minus the final
kinetic energy (positive when energy is lost). Take $u_2=0$ to keep the algebra short; then
$v=m_1u_1/(m_1+m_2)$, and $\tfrac12(m_1+m_2)v^2=\tfrac12 m_1u_1^2\cdot\dfrac{m_1}{m_1+m_2}$, so
\[
K_{\text{lost}}=\tfrac12 m_1u_1^2\Big(1-\frac{m_1}{m_1+m_2}\Big)
=\tfrac12\,\frac{m_1 m_2}{m_1+m_2}\,u_1^2.
\]
The lost fraction is $\dfrac{K_{\text{lost}}}{K_i}=\dfrac{m_2}{m_1+m_2}$.
\stepn{3} \textbf{Read the limits.} If a light body hits a massive one ($m_2\gg m_1$), nearly \emph{all}
the kinetic energy is lost (a bullet into a massive block). If a massive body hits a light one
($m_1\gg m_2$), almost none is lost; the heavy body barely notices. The combination
$m_1m_2/(m_1+m_2)$ is the \textbf{reduced mass} (a mass, in \si{kg}, smaller than either); it appears
whenever the motion of two interacting bodies is split into a centre-of-mass (CM) part and a relative
part, as in the next step.
\stepn{4} \textbf{Why sticking loses the most.} At any instant, before or after the collision, let the
bodies have velocities $v_1,v_2$. Write each as the CM velocity $V=\dfrac{m_1v_1+m_2v_2}{m_1+m_2}$ plus a
share of the relative velocity $w=v_1-v_2$ (a signed scalar, in \si{m/s}):
$v_1=V+\dfrac{m_2}{m_1+m_2}\,w$ and $v_2=V-\dfrac{m_1}{m_1+m_2}\,w$. Substituting into
$\tfrac12m_1v_1^2+\tfrac12m_2v_2^2$, the cross terms, $\pm m_1m_2Vw/(m_1+m_2)$, cancel and
\[
K=\tfrac12(m_1+m_2)V^2+\tfrac12\,\frac{m_1m_2}{m_1+m_2}\,w^2.
\]
Momentum fixes $V$, so the first term is the same before and after any collision. The loss is therefore
greatest when the relative velocity afterwards is zero, that is, when the bodies stick, and it then equals
$\tfrac12\dfrac{m_1m_2}{m_1+m_2}(u_1-u_2)^2$ for any initial velocities (with $u_2=0$ this is Step 2).
\end{derivbox}

\begin{trapbox}{sticking does not lose all the kinetic energy}
Only the kinetic energy of \emph{relative} motion can be lost; the centre-of-mass term survives every
collision. All of it disappears only when $V=0$, as when two equal lumps of putty meet head-on at equal
speeds.
\end{trapbox}


\begin{exbox}{Railway coupling}
A \SI{4000}{kg} wagon moving at \SI{3}{m/s} couples onto a stationary \SI{2000}{kg} wagon. Find the
common speed and the fraction of kinetic energy lost.
\stepn{1} \textbf{Momentum.} $v=\dfrac{4000\times3+0}{4000+2000}=\dfrac{12000}{6000}=\SI{2}{m/s}$.
\stepn{2} \textbf{Fraction lost.} $\dfrac{K_{\text{lost}}}{K_i}=\dfrac{m_2}{m_1+m_2}=\dfrac{2000}{6000}
=\tfrac13\approx\SI{33}{\percent}$.
\stepn{3} \textbf{Check directly.} $K_i=\tfrac12(4000)(3^2)=\SI{18000}{J}$;
$K_f=\tfrac12(6000)(2^2)=\SI{12000}{J}$; lost $=\SI{6000}{J}$, which is indeed $\tfrac13$ of $K_i$.
\ans{\SI{2}{m/s}; $\tfrac13$ lost}
\end{exbox}

\section{Partially Elastic Collisions}\label{sec:partial12}

Real collisions lie between the two extremes. Kinetic energy is no longer conserved, so the coefficient
of restitution supplies the second equation instead, and the kinetic-energy split just derived gives the
loss.

\begin{derivbox}{Final velocities and energy lost for a given $e$}
Masses $m_1,m_2$ with velocities $u_1,u_2$ collide head-on with coefficient of restitution $e$.

\stepn{1} \textbf{Two equations.} Momentum and the restitution law give
\[
m_1u_1+m_2u_2=m_1v_1+m_2v_2,\qquad v_2-v_1=e\,(u_1-u_2).
\]
\stepn{2} \textbf{Solve.} Substituting $v_2=v_1+e(u_1-u_2)$ into the momentum equation and collecting
terms,
\[
v_1=\frac{(m_1-em_2)\,u_1+(1+e)\,m_2u_2}{m_1+m_2},\qquad
v_2=\frac{(m_2-em_1)\,u_2+(1+e)\,m_1u_1}{m_1+m_2}.
\]
Setting $e=1$ recovers the elastic formulas; setting $e=0$ makes $v_1=v_2$, the common velocity of
\S\ref{sec:inelastic12}.
\stepn{3} \textbf{Energy lost.} In the split of the kinetic energy derived in \S\ref{sec:inelastic12}, the
centre-of-mass term $\tfrac12(m_1+m_2)V^2$ does not change, while the relative velocity $w$ goes from $u_1-u_2$ to $-e(u_1-u_2)$. Hence
\[
\boxed{\;K_{\text{lost}}=\tfrac12\,\frac{m_1m_2}{m_1+m_2}\,(1-e^2)\,(u_1-u_2)^2\;}
\]
This is zero for $e=1$ and greatest for $e=0$: a fraction $1-e^2$ of the kinetic energy of
\emph{relative} motion is lost.
\end{derivbox}

\textbf{Collision with a fixed surface.} An important special case is a ball hitting the floor. When
body~2 is a floor or wall attached to the Earth, it does not move measurably, and the restitution law
alone gives the rebound: with $u_2=v_2=0$, $v_1=-e\,u_1$. (The ball's momentum is \emph{not} conserved;
the floor supplies the impulse that reverses it. The momentum of ball plus Earth is conserved, but with
$K=p^2/2M$ and the Earth's mass $M$ enormous, the Earth's share of kinetic energy is negligible.) A ball
dropped from height $h$ reaches the floor at $\sqrt{2gh}$ (from $mgh=\tfrac12mv^2$) and leaves at
$e\sqrt{2gh}$, so it rises to $e^2h$; after $n$ bounces the height is $e^{2n}h$. Measuring the first
rebound height is therefore the simplest way to find $e$: $e=\sqrt{h_1/h}$.

\section{Collisions in Two Dimensions}

Most real collisions are glancing, not head-on: the velocities are not along the line of impact, so the
bodies fly off in new directions. Momentum is still conserved, but now as a \emph{vector} equation, so
\emph{component by component}. Counting shows what else is needed: in a plane the two final velocities
have four components; momentum gives two equations and, for an elastic collision, energy a third. One
more datum, usually a measured angle, must be supplied, which is why such problems always state one.

\begin{figure}[htbp]
\centering
\tikzfig{collision_oblique}
\figcaption{fig:colloblique}{A glancing collision of two smooth spheres. Resolve the incoming velocity along the line of centres ($u\cos\phi$) and perpendicular to it ($u\sin\phi$). The smooth contact force acts only along the line of centres, so only that component takes part in the collision; for equal masses colliding elastically it is handed over entirely, and the two balls leave at right angles.}
\end{figure}

\begin{methodbox}{Oblique collisions: conserve each component}
Choose axes (usually along and perpendicular to the initial motion), and write $u_{1x}$ for the
$x$-component of body~1's initial velocity, and so on. Momentum conservation gives one equation per axis:
\[
m_1u_{1x}+m_2u_{2x}=m_1v_{1x}+m_2v_{2x},\qquad m_1u_{1y}+m_2u_{2y}=m_1v_{1y}+m_2v_{2y}.
\]
If the collision is also elastic, add $\tfrac12 m_1u_1^2+\tfrac12 m_2u_2^2=\tfrac12 m_1v_1^2+\tfrac12
m_2v_2^2$, where now $u_1$, $v_1$, \dots\ are speeds (one scalar equation). If instead the collision has a coefficient of restitution $e<1$ and the bodies are smooth
spheres, apply $e$ only to the velocity components along the line of centres; the components perpendicular
to it are unchanged, because a smooth contact force has no perpendicular component. A common special case: when equal masses collide elastically and one is initially at rest,
they fly apart at \textbf{$90^\circ$} to each other (Figure~\ref{fig:colloblique}; unless the hit is head-on, when the mover stops).
Proof: with $\vc{u}$ the incoming velocity and the mass cancelled, momentum gives $\vc{u}=\vv_1+\vv_2$ and
energy gives $u^2=v_1^2+v_2^2$.
Squaring the first, $u^2=v_1^2+v_2^2+2\,\vv_1\cdot\vv_2$, and subtracting the second leaves
$\vv_1\cdot\vv_2=0$.
\end{methodbox}

\begin{exbox}{A glancing hit}
A ball moving at \SI{5}{m/s} along $x$ strikes an identical stationary ball. After the elastic collision
the incoming ball moves at $30^\circ$ above the $x$-axis. Show the second ball moves at $60^\circ$ below,
and find both speeds.
\stepn{1} \textbf{Equal-mass elastic $\Rightarrow$ $90^\circ$ apart} (the special case in the box). So the
second ball moves at $90^\circ$ to the first, on the other side of the $x$-axis (its $y$-momentum must
cancel the first ball's): $30^\circ-90^\circ=-60^\circ$, that is, $60^\circ$ below the $x$-axis.
\stepn{2} \textbf{Conserve momentum per axis} to get the two speeds $v_1$, $v_2$ (mass cancels):
\[
x:\ 5=v_1\cos30^\circ+v_2\cos60^\circ,\qquad
y:\ 0=v_1\sin30^\circ-v_2\sin60^\circ.
\]
\stepn{3} \textbf{Solve.} From $y$: $v_1(0.5)=v_2(0.866)$, so $v_1=1.732\,v_2$. Substitute into $x$:
$5=1.732v_2(0.866)+v_2(0.5)=2.0\,v_2$, giving $v_2=\SI{2.5}{m/s}$ and $v_1=\SI{4.33}{m/s}$.
\stepn{4} \textbf{Check energy.} $v_1^2+v_2^2=18.75+6.25=25=5^2$ \checkmark, kinetic energy is
conserved, as it must be for an elastic collision. \ans{$v_1=\SI{4.33}{m/s}$, $v_2=\SI{2.5}{m/s}$, at $90^\circ$}
\end{exbox}

\section{Summary}\label{sec:summary12}
Momentum is conserved in every collision between free bodies, along each axis separately; kinetic energy
only in an elastic one, and the coefficient of restitution measures where a collision sits in between.

\begin{itemize}
\item $e=\dfrac{v_2-v_1}{u_1-u_2}$: $e=1$ elastic, $e=0$ perfectly inelastic, $0<e<1$ the real case.
\item General head-on collision: momentum plus $v_2-v_1=e(u_1-u_2)$; the kinetic energy lost is
      $\tfrac12\dfrac{m_1m_2}{m_1+m_2}(1-e^2)(u_1-u_2)^2$.
\item Elastic, target at rest: equal masses exchange velocities; a heavy body hitting a light one sends it off
      at about $2u_1$; a light body hitting a heavy one rebounds at about $-u_1$.
\item Perfectly inelastic: common velocity $v=\dfrac{m_1u_1+m_2u_2}{m_1+m_2}$, and this case loses the most
      kinetic energy consistent with momentum conservation.
\item A ball dropped from height $h$ onto a fixed floor with restitution $e$ rebounds to $e^2h$; the ball's
      momentum is not conserved, because the floor supplies an external impulse.
\end{itemize}

\section{Problem Types}\label{sec:types12}
Five types cover nearly all collision questions, and every one begins with the same question:
\emph{which quantity is conserved here}? Answer that first and the rest is algebra.

\archq{Perfectly inelastic: find the common velocity}
{``Two bodies collide and stick (or a bullet embeds), find the final velocity, or the energy lost.''}
{Momentum only: $v=\dfrac{m_1u_1+m_2u_2}{m_1+m_2}$; then compute the energy lost,
$\tfrac12\dfrac{m_1m_2}{m_1+m_2}(u_1-u_2)^2$.}
{Using kinetic-energy conservation: the sticking \emph{is} the energy loss.}
{A \SI{20}{g} bullet at \SI{400}{m/s} embeds in a \SI{1.98}{kg} block at rest. Their common speed is:}
{\SI{4}{m/s}}
{\SI{2}{m/s}}
{\SI{8}{m/s}}
{\SI{40}{m/s}\archans{A}{$v=\dfrac{0.02\times400}{0.02+1.98}=\dfrac{8}{2}=\SI{4}{m/s}$. Option (C) is
the momentum, \SI{8}{kg.m/s}, not yet divided by the total mass.}}

\archq{Elastic, equal masses: the velocity swap}
{``An elastic head-on collision between equal masses, find the final velocities.''}
{Equal masses in 1-D elastic collision simply \emph{exchange} velocities: $v_1=u_2$, $v_2=u_1$.
If the target was at rest, the incoming body stops dead.}
{This exact swap holds \emph{only} for equal masses and an elastic collision; do not apply it
otherwise.}
{A ball moving at \SI{6}{m/s} elastically strikes an identical ball at rest head-on. Afterwards:}
{both move at \SI{3}{m/s}}
{first stops, second moves at \SI{6}{m/s}}
{first rebounds at \SI{6}{m/s}}
{both move at \SI{6}{m/s}\archans{B}{Equal-mass elastic collision swaps velocities: the first stops
($v_1=0$) and the second leaves with \SI{6}{m/s}. Option (A) conserves momentum but is the sticking
(perfectly inelastic) result, which loses half the kinetic energy.}}

\archq{Coefficient of restitution: the bouncing ball}
{``A ball dropped from $h_0$ rebounds to $h_1$, find $e$,'' or ``find the height after $n$ bounces,''
or ``find the total distance or time.''}
{Speeds relate to heights by $v=\sqrt{2gh}$, so
$\boxed{e=\sqrt{h_1/h_0}}$ and after $n$ bounces $h_n=h_0e^{2n}$. Total distance
$=h_0\dfrac{1+e^2}{1-e^2}$; total time $=\sqrt{\dfrac{2h_0}{g}}\cdot\dfrac{1+e}{1-e}$.}
{Writing $e=h_1/h_0$. The coefficient of restitution is a ratio of \emph{speeds}, and heights go as the
\emph{square} of speed, hence the square root.}
{A ball dropped from \SI{10}{m} rebounds to \SI{2.5}{m}. Its coefficient of restitution is:}
{$1$}{$0.75$}{$0.5$}{$0.25$\archans{C}{$e=\sqrt{h_1/h_0}=\sqrt{2.5/10}=\sqrt{0.25}=0.5$. Option
(D) is the trap of taking the ratio of heights directly.}}

\archq{Fraction of energy lost in a perfectly inelastic collision}
{``Two bodies stick together, what fraction of the kinetic energy is lost?''}
{Momentum gives $v=\dfrac{m_1u_1}{m_1+m_2}$ (target at rest), and the lost fraction is
$\boxed{\dfrac{K_{\text{lost}}}{K_i}=\dfrac{m_2}{m_1+m_2}}$. So a light body hitting a heavy one loses nearly
everything; a heavy hitting a light one loses almost nothing.}
{Assuming half the energy is always lost, or using the striker's mass on top. The fraction depends on
the mass ratio, with the \emph{target's} mass in the numerator.}
{A \SI{2}{kg} body at \SI{3}{m/s} strikes a stationary \SI{4}{kg} body and they stick. The fraction of
kinetic energy lost is:}
{$\tfrac13$}{$\tfrac12$}{$\tfrac23$}{$\tfrac34$\archans{C}{Fraction lost
$=\dfrac{m_2}{m_1+m_2}=\dfrac46=\dfrac23$. (Check: $K_i=\SI{9}{J}$, $v=\SI{1}{m/s}$,
$K_f=\SI{3}{J}$.) Option (A) puts the striker's mass on top; (B) is the ``always half'' guess.}}

\archq{Partially elastic collision using $e$}
{``Two bodies collide with coefficient of restitution $e$, find the final velocities.''}
{Use momentum conservation \emph{together with} the restitution equation
$e=\dfrac{v_2-v_1}{u_1-u_2}$ (\S\ref{sec:partial12}). For equal masses with the target at rest this gives
$v_1=\dfrac{(1-e)u}{2}$ and $v_2=\dfrac{(1+e)u}{2}$. Setting $e=1$ recovers the velocity swap;
$e=0$ gives the common velocity $u/2$.}
{Using energy conservation when $e<1$. Energy is \emph{not} conserved; the restitution equation
replaces it.}
{Two identical balls collide head-on, one at $u$ and one at rest, with $e=0.5$. Their final velocities
are:}
{$0$ and $u$}{$3u/4$ and $u/4$}{$u/2$ and $u/2$}{$u/4$ and $3u/4$\archans{D}{
$v_1=\dfrac{(1-0.5)u}{2}=\dfrac u4$ and $v_2=\dfrac{(1+0.5)u}{2}=\dfrac{3u}{4}$. Momentum checks:
$\tfrac u4+\tfrac{3u}4=u$.\ \checkmark\ Option (B) swaps them, which would have the striker pass
through the ball ahead of it.}}


\section{Practice Problems}\label{sec:practice12}

The twenty problems below cover the whole of Part C, \chsref{ch:energy}{ch:collisions}. Each chains at
least two ideas, as NEET and JEE questions usually do; the numbers come out clean if the physics is
right. Take $g=\SI{10}{m/s^2}$ throughout, and attempt each fully before reading the worked answers.

\begin{enumerate}[label=\textbf{\arabic*.},itemsep=4pt]
\item \textbf{[Vector work $+$ work--energy theorem]} A force $\vF=(4\uhat{\imath}+3\uhat{\jmath})\,\si{N}$ acts
on a \SI{2}{kg} body as it moves through displacement $\vc{d}=(4\uhat{\imath}-3\uhat{\jmath})\,\si{m}$; friction
simultaneously does $\SI{-10}{J}$ of work. If the body started at \SI{2}{m/s}, find its final speed.

\item \textbf{[Variable force]} A force $F(x)=(12x-3x^2)\,\si{N}$ ($x$ in metres) acts along $x$ on a
\SI{4}{kg} body, which passes $x=0$ moving at \SI{3}{m/s} in the $+x$ direction. Find its speed at
$x=\SI{4}{m}$.

\item \textbf{[Average power, two energy sinks]} A pump raises \SI{200}{kg} of water per minute through
\SI{10}{m} and ejects it at \SI{5}{m/s}. Find the power delivered to the water.

\item \textbf{[Power--speed trade-off]} A \SI{1000}{kg} car has maximum engine power \SI{60}{kW} against a
constant total resistance of \SI{1500}{N}. Find (a) its top speed, and (b) its acceleration at
\SI{20}{m/s} with the engine at full power.

\item \textbf{[Potential-energy landscape]} A particle moves in $U(x)=x^3-6x^2+9x$ (joules, $x$ in
metres). Locate the equilibria, classify their stability, and find the force at $x=0$.

\item \textbf{[Energy $+$ circular dynamics]} A block released from height $h$ on a smooth track must
complete a vertical loop of radius $R$. Find the minimum $h$, and the normal force on the block at the
\emph{bottom} of the loop when released from $h=3R$.

\item \textbf{[Gravity feeding a spring]} A \SI{2}{kg} block slides from rest \SI{2.0}{m} down a smooth
$30^\circ$ incline into a spring of stiffness \SI{200}{N/m} aligned with the incline. Find the maximum
compression.

\item \textbf{[Round trip on a rough incline]} A block projected up a $45^\circ$ rough incline returns to
its start with exactly half its launch speed. Find $\mu$.

\item \textbf{[System energy with friction]} A \SI{2}{kg} block on a table ($\mu=0.5$) is connected by a
light inextensible string over a smooth light pulley to a hanging \SI{3}{kg} mass. Released from rest, find the common speed after the hanging
mass falls \SI{1}{m}.

\item \textbf{[Instantaneous power from kinematics]} A \SI{1}{kg} particle moves as $x(t)=t^3/3$ (SI).
Find the power delivered to it at $t=\SI{2}{s}$.

\item \textbf{[Elastic collision, unknown mass]} A ball of mass $m$ moving at speed $u$ strikes a
stationary mass $M$ head-on and elastically, rebounding at $u/3$. Find $M$ and the fraction of the
initial kinetic energy transferred to $M$.

\item \textbf{[Perfectly inelastic in 2-D]} A \SI{2}{kg} body moving east at \SI{6}{m/s} collides and
sticks with a \SI{4}{kg} body moving north at \SI{3}{m/s}. Find the final velocity (magnitude and
direction) and the kinetic energy lost.

\item \textbf{[Ballistic block on a rough floor]} A \SI{20}{g} bullet at \SI{500}{m/s} embeds in a
\SI{1.98}{kg} block resting on a floor with $\mu=0.4$. How far does the block slide?

\item \textbf{[Restitution: heights and total time]} A ball is dropped from \SI{1.25}{m} onto a floor
with $e=0.8$. Find the rebound height after the \emph{second} bounce, and the total time until bouncing
ceases.

\item \textbf{[A chain of different collisions]} Ball A ($m$, speed $u$) strikes identical ball B (at
rest) elastically; B then strikes C (mass $2m$, at rest) and sticks. What fraction of the original
kinetic energy survives at the end?

\item \textbf{[Symmetric partial restitution]} Two identical balls approach head-on, each at speed $u$,
and collide with $e=0.5$. Find their final velocities and the fraction of kinetic energy lost.

\item \textbf{[Pendulum--collision--friction, three stages]} A \SI{1}{kg} bob on a string of length
$\ell=\SI{1}{m}$ is
released from rest with the string horizontal. At the lowest point it strikes a \SI{1}{kg} block head-on
and elastically; the block then
slides on a floor with $\mu=0.5$. How far does the block slide, and what does the bob do afterwards?

\item \textbf{[Spring pushing two masses apart]} A compressed spring storing \SI{12}{J} is released
between a \SI{1}{kg} and a \SI{3}{kg} block on a smooth floor. Find both final speeds and each block's
kinetic energy. Which block gets more, and why?

\item \textbf{[Impulse at a wall]} A \SI{0.2}{kg} ball hits a wall normally at \SI{25}{m/s} and rebounds
with $e=0.8$. If contact lasts \SI{10}{ms}, find the impulse on the wall and the average force.

\item \textbf{[Scaling with $P\propto v^3$]} At top speed an engine's power is spent entirely against
aerodynamic drag, which grows as the square of the speed ($F=cv^2$, with $c$ a constant). If the engine
power is doubled, by what percentage does the top speed rise?
\end{enumerate}

\subsection*{Answers}
\begin{enumerate}[label=\textbf{\arabic*.},itemsep=5pt]
\item Two forces do work, so add their works to get the net work: $W_{\vF}=\vF\cdot\vc{d}=(4)(4)+(3)(-3)=\SI{7}{J}$;
$W_{\text{net}}=7-10=\SI{-3}{J}$.
$K_i=\tfrac12(2)(2^2)=\SI{4}{J}$, so $K_f=\SI{1}{J}$ and $v_f=\sqrt{2K_f/m}=\SI{1}{m/s}$. \ans{\SI{1}{m/s}}

\item The force varies, so integrate for the work, then apply the work--energy theorem:
$W=\int_0^4(12x-3x^2)\,dx=\big[6x^2-x^3\big]_0^4=96-64=\SI{32}{J}$. Then
$\tfrac12(4)v^2=\tfrac12(4)(3^2)+32=\SI{50}{J}$, so $v=\SI{5}{m/s}$. (The body must be moving at $x=0$:
there $F=0$, and a body at rest at $x=0$ would stay there.) \ans{\SI{5}{m/s}}

\item Each kilogram of water gains $gh$ of potential and $\tfrac12v^2$ of kinetic energy; the mass pumped
per second is $\dot m=200/60\,$kg/s. So $P=\dot m\big(gh+\tfrac12v^2\big)=\tfrac{200}{60}(100+12.5)
=\SI{375}{W}$. \ans{\SI{375}{W}}

\item (a) At top speed, drive force $=$ resistance: $v_{\max}=P/F=60000/1500=\SI{40}{m/s}$.
(b) At \SI{20}{m/s} the engine supplies $F=P/v=\SI{3000}{N}$; net $=\SI{1500}{N}$;
$a=\SI{1.5}{m/s^2}$. \ans{\SI{40}{m/s}; \SI{1.5}{m/s^2}}

\item $F=-\dfrac{dU}{dx}=-(3x^2-12x+9)=-3(x-1)(x-3)$. Equilibria $x=1,3$. Writing $U''$ for
$d^2U/dx^2$, $U''=6x-12$: at $x=1$,
$U''<0$ (maximum, \emph{unstable}); at $x=3$, $U''>0$ (minimum, \emph{stable}). At $x=0$:
$F=\SI{-9}{N}$ (towards $-x$). \ans{$x=1$ unstable, $x=3$ stable; $F(0)=\SI{-9}{N}$}

\item At the top the normal force can only push down, so the centripetal force must be at least $mg$:
$v_{\text{top}}^2\ge gR$ (\chref{ch:circdyn}). Energy from $h$ to the top, divided by $m$:
$g h=\tfrac12 v_{\text{top}}^2+g(2R)\Rightarrow h_{\min}=\tfrac52 R$. For $h=3R$: at the bottom
$v^2=2g(3R)=6gR$, and Newton's law on the circle gives $N-mg=\dfrac{mv^2}{R}=6mg$, so $N=7mg$.
\ans{$h_{\min}=2.5R$; $N=7mg$}

\item The spring force varies, so use energy. At maximum compression $x$ the block is momentarily at rest,
having descended $d+x$ along the incline ($d=\SI{2.0}{m}$), so gravitational energy lost equals spring
energy stored:
$mg\sin30^\circ\,(d+x)=\tfrac12kx^2$. With numbers: $2(10)(0.5)(2.0+x)=100x^2\Rightarrow
10x^2-x-2=0$, so $x=\dfrac{1+\sqrt{1+80}}{20}=\dfrac{1+9}{20}=\SI{0.5}{m}$ (the other root,
$-0.4$, is unphysical). Check: $10(2.5)=25=100(0.5)^2$.\ \checkmark \ans{\SI{0.5}{m}}

\item Let $u$ be the launch speed, $v=u/2$ the return speed, $s$ the distance up the slope and
$\theta=45^\circ$. Friction $\mu mg\cos\theta$ opposes the motion both ways. Energy per unit mass, up:
$\tfrac12u^2=gs(\sin\theta+\mu\cos\theta)$; down: $\tfrac12v^2=gs(\sin\theta-\mu\cos\theta)$. Dividing
removes the unknown $s$.
Given $v=u/2$: $\dfrac{\sin\theta-\mu\cos\theta}{\sin\theta+\mu\cos\theta}=\dfrac14\Rightarrow
\mu=\tfrac35\tan\theta=\tfrac35$ at $45^\circ$. \ans{$\mu=0.6$}

\item Treat block ($m_1=\SI{2}{kg}$) and hanging mass ($m_2=\SI{3}{kg}$) as one system, so the tension
does no net work. As $m_2$ falls $h=\SI{1}{m}$, gravity supplies $m_2gh=\SI{30}{J}$ and friction over the
block's equal slide costs $\mu m_1gh=\SI{10}{J}$; the rest is kinetic:
$\tfrac12(m_1+m_2)v^2=20\Rightarrow v^2=8$, $v=2\sqrt2\approx\SI{2.83}{m/s}$.
\ans{$2\sqrt{2}\approx\SI{2.83}{m/s}$}

\item $P=Fv$ needs both force and velocity, so differentiate: $v=dx/dt=t^2$, $a=dv/dt=2t$, $F=ma=2t$,
so $P=Fv=2t^3$. At $t=2$: $P=\SI{16}{W}$. \ans{\SI{16}{W}}

\item Use the elastic formula for a target at rest, with the rebound velocity $-u/3$ (negative: it
reverses): $v_1=\dfrac{m-M}{m+M}u=-\dfrac{u}{3}\Rightarrow 3(M-m)=m+M\Rightarrow M=2m$. Then
$v_2=\dfrac{2m}{m+M}u=\dfrac{2u}{3}$ and the transferred fraction is
$\dfrac{\tfrac12(2m)(2u/3)^2}{\tfrac12mu^2}=\dfrac89$. \ans{$M=2m$; $8/9$}

\item Sticking means momentum only, as a vector. With $\uhat{\imath}$ east and $\uhat{\jmath}$ north,
the total momentum is $\vc P=(12\,\uhat{\imath}+12\,\uhat{\jmath})\,\si{kg.m/s}$; $v=\dfrac{|\vc P|}{6}=
\dfrac{12\sqrt2}{6}=2\sqrt2\approx\SI{2.83}{m/s}$ at $45^\circ$ north of east.
$K_i=36+18=\SI{54}{J}$, $K_f=\tfrac12(6)(8)=\SI{24}{J}$: lost \SI{30}{J}. \emph{Momentum is conserved
per component; energy is not.} \ans{$2\sqrt2\,\si{m/s}$ at $45^\circ$ north of east; \SI{30}{J} lost}

\item Stage 1 (collision, momentum only): $v=\dfrac{0.02\times500}{2.0}=\SI{5}{m/s}$. Stage 2 (sliding,
work--energy): friction $\mu Mg$ on the combined mass $M=\SI{2.0}{kg}$ removes $\tfrac12Mv^2$, so
$d=\dfrac{v^2}{2\mu g}=\dfrac{25}{8}=\SI{3.125}{m}$. \ans{\SI{3.125}{m}}

\item Heights scale by $e^2$ per bounce: after two, $h_2=h_0e^4=1.25(0.8)^4=\SI{0.512}{m}$. Total time:
first fall $t_0=\sqrt{2h_0/g}=\SI{0.5}{s}$; each later up-and-down trip lasts $e$ times the one before,
so the times form a geometric series summing to $T=t_0\dfrac{1+e}{1-e}=0.5\times9=\SI{4.5}{s}$ (the
bouncing-ball problem type). \emph{Infinitely many bounces, finite time.}
\ans{\SI{0.512}{m}; \SI{4.5}{s}}

\item A--B elastic, equal masses: swap, so A stops and B moves at $u$. B--C sticky:
$v=\dfrac{mu}{3m}=\dfrac u3$. Surviving fraction $=\dfrac{\tfrac12(3m)(u/3)^2}{\tfrac12mu^2}=\dfrac13$.
\ans{$\tfrac13$}

\item The total momentum is zero ($mu-mu$), so they leave with equal and opposite velocities, of size
$v$. Separation
$=e\times$approach: $2v=0.5(2u)\Rightarrow v=u/2$ each, reversed. Energy lost
$=1-e^2=\SI{75}{\percent}$ (all the kinetic energy here is relative motion, since the centre of mass is
at rest). \ans{$u/2$ each, reversed; \SI{75}{\percent}}

\item Stage 1 (energy): the bob falls through $\ell$, so $v=\sqrt{2g\ell}=\sqrt{20}=2\sqrt5\,\si{m/s}$ at
the bottom. Stage 2 (equal-mass
elastic): velocities swap, the bob \emph{stops dead}, the block leaves at $2\sqrt5$. Stage 3
(work--energy): $d=\dfrac{v^2}{2\mu g}=\dfrac{20}{10}=\SI{2}{m}$. \ans{\SI{2}{m}; the bob hangs at rest}

\item The spring is internal, so momentum stays zero: $(1)v_1=(3)v_2$, that is, $v_1=3v_2$ (in opposite
directions). All the stored energy becomes kinetic: $\tfrac12(1)v_1^2+\tfrac12(3)v_2^2=12\Rightarrow6v_2^2=12$, so
$v_2=\sqrt2\,\si{m/s}$, $v_1=3\sqrt2\,\si{m/s}$; $K_1=\SI{9}{J}$, $K_2=\SI{3}{J}$. The \emph{lighter} block gets more:
equal momenta mean $K=p^2/2m$, inversely proportional to mass. \ans{$3\sqrt2$ and $\sqrt2\,\si{m/s}$; \SI{9}{J} and \SI{3}{J}}

\item Fixed wall, so the rebound speed is $eu=\SI{20}{m/s}$. Impulse is change of momentum, and the
velocity reverses, so the speeds add: $J=0.2(25+20)=\SI{9}{N.s}$ (the wall receives the same, by the third
law). Average force $\bar F=J/\Delta t=9/0.01=\SI{900}{N}$.
\ans{\SI{9}{N.s}; \SI{900}{N}}

\item At top speed $P=Fv=(cv^2)v=cv^3$, so $v\propto P^{1/3}$. Doubling $P$ multiplies $v$ by
$2^{1/3}\approx1.26$: about a \SI{26}{\percent} rise. \emph{Eight times the power for double the
speed, read backwards.} \ans{$\approx\SI{26}{\percent}$}
\end{enumerate}

\section{Challenge Problems}\label{sec:challenge12}
The three problems below join the methods of \chsref{ch:energy}{ch:collisions}, as harder examinations
do: a collision governed by momentum next to a motion governed by energy, or an energy argument with a
hidden subtlety. Finding the boundary between the stages is most of the work. Take
$g=\SI{9.8}{m/s^2}$.

\begin{challengebox}{Spring, slide, and stop (builds on energy methods)}
A spring of stiffness $k=\SI{200}{N/m}$ is compressed by $x_0=\SI{0.10}{m}$ and released, launching a
\SI{0.50}{kg} block (not attached to the spring) along a rough horizontal floor with $\mu=0.3$. (a) How far
does the block travel, from its release point, before stopping? (b) Where is the block moving fastest,
and at what speed?
\challengeans{(a) The spring force varies and the block ends at rest, so energy accounting from start to
finish avoids the motion in between: all the spring's stored energy is eventually removed by friction,
which acts from the moment of release with $f=\mu mg=0.3\times0.5\times9.8=\SI{1.47}{N}$.
\[
\tfrac12 k x_0^2=f\,d\ \Rightarrow\ d=\frac{\tfrac12(200)(0.10)^2}{1.47}
=\frac{1.0}{1.47}\approx\SI{0.68}{m}.
\]
The block does start and does leave the spring: the initial spring force, \SI{20}{N}, exceeds $f$, and
at the natural length the block still has $1.0-1.47\times0.10\approx\SI{0.85}{J}$ of kinetic energy.

(b) The tempting answer, ``at the moment it leaves the spring,'' is wrong. The block speeds up only while
the spring force exceeds friction, $kx>f$, so the maximum speed occurs at compression
$x_1=f/k=1.47/200=\SI{7.35}{mm}$, slightly \emph{before} the natural length. Between release and $x_1$,
\[
\tfrac12mv_{\max}^2=\tfrac12k\big(x_0^2-x_1^2\big)-f\,(x_0-x_1)\approx0.9946-0.1362=\SI{0.858}{J},
\]
so $v_{\max}\approx\SI{1.853}{m/s}$, against \SI{1.847}{m/s} at the natural length. The difference is
small here, but the principle (maximum speed where the net force vanishes) is what examiners test.}
\end{challengebox}

\begin{challengebox}{A bullet through a pendulum bob (builds on collisions, energy and circular motion)}
A bullet of mass $m=\SI{20}{g}$ travelling horizontally at $u=\SI{350}{m/s}$ passes straight through a
block of mass $M=\SI{1.0}{kg}$ hanging at rest on a light string of length $\ell=\SI{0.70}{m}$, and
emerges at $u'=\SI{105}{m/s}$. Treat the block as a particle and neglect the mass the bullet loses. (a)
Find the block's speed $v_B$ just after the bullet leaves. (b) Does the block complete a vertical circle?
If not, find the angle at which the string goes slack. (c) What fraction of the bullet's kinetic energy
is dissipated in the block?
\challengeans{\textbf{(a) Collision: momentum only.} The bullet passes through in a very short time, so
the block barely moves, the string stays vertical and exerts no horizontal impulse; horizontal momentum is
conserved. Kinetic energy is not, because the bullet tears its way through.
\[
mu=mu'+Mv_B\ \Rightarrow\ v_B=\frac{m(u-u')}{M}=\frac{0.020\times(350-105)}{1.0}=\SI{4.9}{m/s}.
\]
\textbf{(b) Swing: energy plus the radial equation.} A particle on a string completes the circle only if
$v_B^2\ge5g\ell$ (\chref{ch:circdyn}), and rises above the pivot's level only if $v_B^2>2g\ell$. Here
$v_B^2=\SI{24.01}{m^2/s^2}$, while $5g\ell=\SI{34.3}{m^2/s^2}$ and $2g\ell=\SI{13.72}{m^2/s^2}$. Since
$2g\ell<v_B^2<5g\ell$, the block rises above the level of the pivot but the string goes slack before the
top. Let $\theta$ be the angle of the string from the downward vertical, $v$ the block's speed there and
$T$ the tension. Energy gives $v^2=v_B^2-2g\ell(1-\cos\theta)$, and the radial equation (tension plus the
inward component of gravity, $-Mg\cos\theta$, supply $Mv^2/\ell$) gives
$T=\dfrac{Mv^2}{\ell}+Mg\cos\theta$. The string goes slack when $T=0$, that is when
$v^2=-g\ell\cos\theta$. Eliminating $v^2$,
\[
\cos\theta=\frac{2g\ell-v_B^2}{3g\ell}=\frac{13.72-24.01}{20.58}=-\frac12\ \Rightarrow\ \theta=120^\circ,
\]
which is $30^\circ$ above the horizontal through the pivot, at a height $\ell(1-\cos\theta)=1.5\,\ell=\SI{1.05}{m}$
above the lowest point. There the speed is $v=\sqrt{g\ell/2}\approx\SI{1.85}{m/s}$, and the block leaves the
circle as a projectile.

\textbf{(c) Energy dissipated.} Before the collision there is $\tfrac12(0.020)(350)^2=\SI{1225}{J}$; just
after it, the bullet keeps $\tfrac12(0.020)(105)^2=\SI{110.25}{J}$ and the block has
$\tfrac12(1.0)(4.9)^2\approx\SI{12.0}{J}$. The difference, about \SI{1103}{J}, is a fraction
$1102.7/1225\approx\SI{90}{\percent}$, which goes into heat and deformation. Only about \SI{1}{\percent}
of the bullet's energy reaches the block as kinetic energy, so the energy method could not have been
used through the collision.}
\end{challengebox}

\begin{challengebox}{Two stacked balls (builds on elastic collisions in a moving frame)}
A small ball of mass $m$ rests on top of a large ball of mass $M=3m$, with a tiny gap between them, and
the pair is dropped from height $h$ onto a hard floor. All collisions are elastic and the balls' sizes
are negligible compared with $h$. How high does the small ball rise, and what does the large ball do?
\challengeans{The trick is to split the event into two separate collisions. Both balls reach the floor
region at speed $v=\sqrt{2gh}$. Because of the gap, the large
ball bounces off the floor first and is moving \emph{up} at $v$ (elastic, fixed floor: rebound with
$e=1$) while the small ball is still moving \emph{down} at $v$. Now apply the two-body elastic formulas,
with the small ball as body~1, upward positive,

$u_{\text{small}}=-v$ and $u_{\text{large}}=+v$:
\[
v_{\text{small}}=\frac{m-M}{m+M}(-v)+\frac{2M}{m+M}(+v)=\frac{3M-m}{M+m}\,v=2v,
\]
\[
v_{\text{large}}=\frac{M-m}{M+m}(+v)+\frac{2m}{M+m}(-v)=\frac{M-3m}{M+m}\,v=0.
\]
The large ball is left at rest on the floor, and the small ball leaves at $2v$, so it rises to
$(2v)^2/2g=4h$. Check: the total energy $(m+3m)gh=4mgh$ has all gone into the small ball,
$mg(4h)=4mgh$. For $M\gg m$ the same formulas give $3v$ and a height of $9h$, which is why a tennis
ball dropped on a basketball shoots so high.}
\end{challengebox}


\section*{Looking Back on Part C}
\addcontentsline{toc}{section}{Looking Back on Part C}

\textbf{Part C introduced the conserved quantities.}
\begin{itemize}
\item \textbf{Work and energy (\chref{ch:energy}).} Work is $\vF\cdot\vc{d}$: only the component along the motion
counts, so a perpendicular force does none. The work--energy theorem $W_{\text{net}}=\Delta K$ holds for
\emph{any} force and trades the whole history of a motion for its endpoints; power $P=\vF\cdot\vv$ is the
rate of that work. For conservative forces we define a potential energy with $F=-dU/dx$, so a minimum of $U$
is a stable equilibrium and a maximum an unstable one. When only conservative forces do work, $K+U$ is
conserved; in general $\Delta(K+U)=W_{\text{nc}}$, and friction from a fixed surface drains mechanical energy
to heat.
\item \textbf{Collisions (\chref{ch:collisions}).} One asymmetry: \emph{momentum is conserved in every
collision between free bodies} (the internal third-law forces cancel), but \emph{kinetic energy only in
an elastic one}. The coefficient of restitution $e$ supplies the second equation for any head-on
collision; a fraction $1-e^2$ of the kinetic energy of relative motion is lost, the most when the bodies
stick.
\end{itemize}
\textbf{The thread into Part D:} energy and momentum were built for particles. Now we give bodies size and
spin: momentum gains a rotational twin, angular momentum, conserved when there is no external torque, and
kinetic energy gains a rotational part.
